\clearpage
\chapterConclusion
\begin{enumerate}
	\item \textbf{Нелінійність задачі:} Результати чітко доводять, що задача нормалізації
		та перенесення колірного простору супутникових даних є глибоко нелінійною.
		Класичні параметричні регресійні моделі (лінійна, поліноміальна) не здатні вловити
		розподіл каналів та призводять до викривлення даних.

	\item \textbf{Перевага нейромереж:} Багатошаровий перцептрон (MLP) завдяки
		наявності прихованих шарів та нелінійних функцій активації найкраще
		адаптується до спектральних характеристик знімків. Результати обробки за допомогою
		MLP візуально демонструють практично ідеальний безшовний перехід між правою та
		лівою частинами результуючого зображення.
\end{enumerate}